\chapter*{Sequences}
\addcontentsline{toc}{chapter}{Sequences}

\section*{Definition and Point of View}

A sequence is simply a function whose domain is the positive integers. We usually write it as
$$\{a_n\}_{n=1}^{\infty} = a_1,a_2,a_3,\dots$$
where $a_n$ is the $n$th term.

The main question is: \emph{what happens to $a_n$ as $n\to\infty$?}

\begin{definition}[Limit of a sequence]
A sequence $\{a_n\}$ converges to $L$ if
$$\lim_{n\to\infty} a_n = L.$$ 
If no finite limit exists, the sequence diverges.
\end{definition}

\section*{Common Examples}

\begin{itemize}
    \item $a_n=\frac{1}{n}$ converges to $0$.
    \item $a_n=\frac{n+1}{n}=1+\frac{1}{n}$ converges to $1$.
    \item $a_n=(-1)^n$ does not settle to one number, so it diverges.
    \item $a_n=n$ grows without bound, so it diverges to infinity.
\end{itemize}

\section*{Limit Laws for Sequences}

If $a_n\to L$ and $b_n\to M$, then the same algebra of limits from ordinary calculus still works:
\begin{align*}
\lim_{n\to\infty}(a_n+b_n) &= L+M, \\
\lim_{n\to\infty}(a_n b_n) &= LM, \\
\lim_{n\to\infty}\frac{a_n}{b_n} &= \frac{L}{M} \qquad (M\neq 0).
\end{align*}

These rules let us treat many sequences exactly like rational functions in a continuous variable.

\section*{The Squeeze Principle}

If
$$a_n \le b_n \le c_n$$
for all sufficiently large $n$, and if
$$\lim_{n\to\infty} a_n = \lim_{n\to\infty} c_n = L,$$
then
$$\lim_{n\to\infty} b_n = L.$$ 

This is especially useful for sequences involving trigonometric expressions, for example
$$-\frac{1}{n} \le \frac{\sin n}{n} \le \frac{1}{n},$$
so $\sin n/n \to 0$.

\section*{Monotone and Bounded Sequences}

A sequence is called \textbf{increasing} if $a_{n+1}\ge a_n$ and \textbf{decreasing} if $a_{n+1}\le a_n$.

\begin{theorem}[Monotone Convergence Theorem]
If a sequence is monotone and bounded, then it converges.
\end{theorem}

This theorem is important because many recursively defined sequences are difficult to analyze directly, but easier to study using monotonicity and bounds.

\section*{A Note on Recursively Defined Sequences}

Some sequences are defined by a rule such as
$$a_{n+1}=f(a_n).$$
If we suspect $a_n\to L$, then taking limits on both sides often gives a fixed-point equation
$$L=f(L).$$
This does not prove convergence by itself, but it helps identify the only possible limit.